% Plan v4/v5 / related work. Compressed 2026-09-08 for the ICLR page budget.
% Every venue is the one in the arXiv primary record; four entries carry no venue there and are
% cited as preprints. Positioning against segal2026provably and cohen2026barriers is in the
% introduction as well as here -- mentioned only in passing, it would read as concealment.

\section{Related work in full}\label{app:related}

\paragraph{Access-free guarantees and their discontents}
Near-access-freeness \citep{vyas2023provable} asks that a served model stay within a divergence of
a safe model that never saw the protected work, and CP-$\Delta$ and CP-$k$ instantiate it by
combining models trained on disjoint shards \citep{vyas2023provable,abad2025cpfuse}. Three recent
results bear directly on ours. \citet{cohen2025blameless} shows that per-query access-freeness does
not compose across queries; \citet{cohen2026barriers} proves barriers for counterfactual credit
attribution, a sibling relaxation of differential privacy. Closest of all,
\citet{segal2026provably} report that CP-style guarantees can be numerically uninformative in
realistic fine-tuning settings and demonstrate extraction against NAF-protected models. We are
downstream of that observation, not upstream of it.

What differs is the shape of the claim, and it is worth being exact about it because the two
results are compatible. Their diagnosis is that memorisation puts sharp probability spikes on a few
tokens, so a divergence taken over \emph{raw} probabilities reads as large even where the models
agree on what to emit; their repair charges relative probabilities instead. That is a change in
\emph{what} is metered. Proposition~\ref{prop:threshold} closes a different lever, the one a
deployer actually has: for every order of the same per-step charge, from KL to max-divergence, the
certificate becomes vacuous at the same budget, so a stronger charge buys precision below the
threshold and never moves it. Neither result subsumes the other --- ours does not say no repair
exists, and theirs does not claim the order is the problem --- but together they leave a deployer
holding a scalar budget with only one lever that moves the threshold, a smaller budget, which
Theorem~\ref{thm:nfl} prices and Section~\ref{sec:orders} measures.
\citet{cohen2026barriers} is the closer structural parallel, and not the one we first read it as:
its non-composition is \emph{autoregressive}, not across queries. Imposing counterfactual credit
attribution on the next-token predictor does not make the sequence-level model satisfy it, and
retrofitting credit onto a model that lacks it needs query complexity exponential in the output
length. Ours is the same shape for a different quantity --- a per-step charge that is honoured at
every step still leaves the sequence-level event $\{\text{output}=x\}$ unprotected once the budget
reaches $S(x)$ --- and it is priced in nats for a single query rather than stated as a barrier.

\paragraph{Other inference-time defences, and the level none of them certifies}
Anchored decoding is one of several ways to intervene at serving time. CP-Fuse combines models
trained on disjoint shards under a balancing property \citep{abad2025cpfuse} and is the second
mechanism we audit; BloomScrub certifies takedowns by Bloom-filter screening with abstention
\citep{zhang2025bloomscrub}; SILO isolates legal risk in a nonparametric datastore
\citep{min2024silo}; others refuse \citep{liu2024shield} or randomise the decoder
\citep{chen2024randomization}, and for diffusion models retrieval carries a NAF guarantee
\citep{golatkar2024cpr} in response to replication in image models \citep{somepalli2023diffusion}.
\citet{alshehyari2026copyshield} benchmark output-, behavioural- and representation-level defences
across two model families; the level that benchmark omits is the certified one, which is exactly
where a low measured leakage and a vacuous certificate can coexist. The closest empirical precedent
for Section~\ref{sec:frontier}'s accounting is VA3 \citep{li2024va3}, which amplifies a
probabilistic copyright guarantee for text-to-image models through repeated interaction; what
composes there is a per-sample probability accumulating over retries, where a budget adds by
construction.

\paragraph{Training-time alternatives, and the evaluation standard they set}
The alternative is to remove memorisation before deployment: drop tokens from the loss
\citep{hans2024goldfish}, optimise directly against regurgitation \citep{chen2025parapo}, or
unlearn the work \citep{eldan2023whos}. Audits of unlearning
\citep{lynch2024eight,shi2025muse,lucki2025adversarial,deeb2024unlearning} recovered the removed
text by relearning, jailbreaks and probing; the word covers objectives that demand different
evaluations \citep{yoon2026position}, and equivalence to retraining is unattainable
\citep{yu2025impossibility}. Part of the appeal of an inference-time certificate is that it
promises something checkable instead --- which is the promise this paper measures, and it inherits
that evaluation standard rather than escaping it.

\paragraph{Measuring memorisation, and what the measurement is a function of}
Section~\ref{sec:onset} is a contribution to this literature as much as it uses it. Extractability
is defined relative to a prefix \citep{carlini2021extracting,carlini2023quantifying}, and how the
probe is built has repeatedly been found to move the number --- so far, in the direction of
\emph{understating} it. A filter that perfectly prevents verbatim memorisation still leaks under
minimally modified style-transfer prompts, so exact-match definitions are too restrictive
\citep{ippolito2023preventing}; and discoverable extraction, which asks whether greedy decoding
reproduces the suffix, ignores the non-determinism of realistic sampling, which is why
\citet{hayes2025measuring} replace the yes-or-no test with a probability over repeated queries.
We met that second effect from the other side: one candidate memoriser scored $0.708$ greedy
against $0.022$ sampled, which is why a pair enters Section~\ref{sec:onset} only on its
\emph{sampled} $k=-1$ recall (Appendix~\ref{app:limitations}). Adversarial prompting moves the
number again \citep{nasr2025scalable}, capacity-based accounts measure a different quantity
entirely \citep{morris2025memorize}, and membership and contamination tests are sensitive to how
the probe is constructed \citep{meeus2024traps,duarte2024decop}. Our finding is one more such
dependence and an unusually mechanical one: because the prefix is specified in \emph{tokens}, two
models with different tokenizers are handed different amounts of the work at the same nominal
setting, and the quantity that moves is not the memoriser but the budget the target's own
surprisal profile demands.

\paragraph{Copyright as a legal question, and which quantity it asks for}
Whether copyright reduces to a privacy-style definition at all is contested
\citep{elkinkoren2024copyright}; memorisation and infringement are distinct questions
\citep{haviv2026separate,cooper2025files}; substantial similarity has been formalised as
description length with and without access to the original \citep{scheffler2022formalizing}; and
whether a model ``contains a copy'' has been argued to turn on how easily the work can be extracted
\citep{lemley2026probabilistic}, which is the quantity Section~\ref{sec:onset} locates. Where the
deployer sits in the liability chain is its own literature
\citep{henderson2023fairuse,lee2024talkin}. We take no position on any of these; we report what a
published $k$ does and does not tell a rights-holder who asks about one work.

\paragraph{Impossibility results for inference-time mitigations}
The template is \citet{zhang2024watermarks}, who prove strong watermarking impossible under stated
oracle assumptions and demonstrate it against the schemes of
\citet{kirchenbauer2023watermark,kuditipudi2024distortionfree} --- watermarking being the other
inference-time intervention whose guarantee is a divergence from the unmodified model, with
distortion-freeness \citep{christ2024undetectable} the case of zero divergence --- and \citet{sadasivan2025can}, who bound any detector's accuracy by the total
variation between the two distributions and observe that the bound goes vacuous as they converge
--- the same proof shape as ours in a neighbouring problem. \citet{kalai2026consensus} give a
Pareto-optimal safe aggregation rule whose robustness notion parallels $k$-NAF at R\'enyi order
infinity, one endpoint of the family we sweep. \citet{feldman2025tradeoffs} derive memorisation
lower bounds from strong data processing inequalities, and \citet{livni2024credit} place credit
attribution between differential privacy and copyright.

\paragraph{What a divergence budget certifies, and how it is metered}
That a KL bound does not imply a worst-case event bound is classical
\citep{sason2016fdivergence,vanerven2014renyi}, and the differential-privacy literature has
developed the conversions and their costs; our contribution is not the inequality but which budgets
survive the requirement that the certificate be informative about a specific work \emph{and} the
mechanism remain usable. The accounting machinery we borrow is the R\'enyi filter of
\citet{feldman2020renyifilter} and fully adaptive composition
\citep{whitehouse2023fully,rogers2016odometers}; the token bucket and its running maximum are
standard objects in queueing and network calculus
\citep{loynes1962stability,cruz1991calculus,leboudec2001network}, which is where our second
boundary comes from.

\paragraph{Decoding-time control and extraction}
Anchored decoding sits in a family of methods that steer a model at inference under a divergence
constraint \citep{mudgal2024controlled,kim2025guaranteed,khalifa2021distributional}.
\citet{huang2025bestofn} characterise inference-time alignment by a coverage coefficient, a
worst-case density ratio rather than an average divergence, and \citet{monteiropaes2026limits} give
the exact reward attainable at a fixed KL budget; neither asks what the budget certifies about a
specific \emph{event}, which is the rights-holder's question. Our attack strategies follow the
extraction literature \citep{carlini2021extracting,nasr2025scalable}, with passages and coverage
estimates from \citet{cooper2025extracting}. Methodologically we follow
\citet{aerni2024evaluations}, who find that evaluations reporting average-case leakage understate
the leakage of the most vulnerable samples by an order of magnitude, and that no empirical defence
they study beats a tuned DP-SGD baseline whose own provable guarantee is vacuous. That is the gap
this paper measures from the other side: between what the budget reports and what the running
maximum controls, and between a certificate being vacuous and a mechanism being useless.

\paragraph{Consensus sampling and DP training, compared directly}
Two neighbours deserve a direct comparison. \citet{kalai2026consensus} aggregate $k$ distributions
and serve only where they agree, abstaining otherwise, with risk competitive with the average of the
safest $s$ and an $R$-robustness bound on leakage and adversarial influence. That is the same
instinct --- buy safety at the aggregation step --- from the opposite premise: they need $k$ models
of which some unknown subset is safe, and they pay for the guarantee by abstaining, where selection
anchoring needs one \emph{named} safe model, never abstains, and prices itself in the same R\'enyi
family the deployed meters charge in, which is what puts the two mechanisms on one axis.
\citet{chen2024randomization} ask whether differentially private \emph{training} can deliver
near-access-freeness and find the cost prohibitive; we take the budget as given at inference and ask
only where it should sit.
